\subsection{EXTERIOR ALGEBRA OF A FREE MODULE}

THEOREM 1. Let M be an A-module with a basis \((e_{\lambda})_{\lambda \in L}\) . Let L be given the structure of a totally ordered set (Set Theory, III, §2, no. 3, Theorem 1) and for every finite subset J of L we write

\[
e _ {J} = e _ {\lambda_ {1}} \wedge e _ {\lambda_ {2}} \wedge \dots \wedge e _ {\lambda_ {n}}\tag{18}
\]

where \((\lambda_k)_{1 \leqslant k \leqslant n}\) is the sequence of elements of \(J\) arranged in increasing order (Set Theory, III, § 5, no. 3, Proposition 6); we write \(e_\varphi = 1\) , the unit element of A. Then the \(e_J\) , where \(J\) runs through the set \(\mathfrak{F}(L)\) of finite subsets of \(L\) , form a basis of the exterior algebra \(\bigwedge(M)\) .

Since the \(e_{\lambda}\) generate the A-module M, every element of \(\bigwedge(\mathbf{M})\) is a linear combination of a finite number of products of the elements \(e_{\lambda}\) and hence (taking account of no. 3, Proposition 5) is a linear combination of a finite number of elements \(e_{J}\) for \(J \in \mathfrak{F}(L)\) . It reduces to proving that the \(e_{J}\) are linearly independent over A. Otherwise, there would exist between these elements a linear relation with coefficients not all zero; the union of the subsets J which correspond to the \(e_{J}\) whose coefficients in this relation are \(\neq 0\) is a finite subset K of L (since there is only a finite number of coefficients \(\neq 0\) ). Let N be the submodule of M generated by the \(e_{\lambda}\) such that \(\lambda \in K\) ; N is a direct factor in M, hence (no. 2) \(\bigwedge(\mathbf{N})\) is identified with a subalgebra of \(\bigwedge(\mathbf{M})\) and, if we show that the \(e_{J}\) with \(J \subset K\) form a basis of \(\bigwedge(\mathbf{N})\) , we shall obtain the desired contradiction.

It therefore reduces to providing Theorem 1 when the basis of M is finite; we may therefore suppose that \(L = [1, m] \subset N\) . For each \(i \in L\) , let \(M_i\) be the free submodule \(Ae_i\) of M; M is the direct sum of the \(M_i\) and \(\bigwedge(M_i)\) is the direct sum of \(\bigwedge^0(M_i) = A\) and \(\bigwedge^1(M_i) = M_i\) (no. 3, Proposition 6). Let \(\bigwedge(M)\) be canonically identified with the A-module the tensor product of the \(\bigwedge(M_i)\) (no. 7, Proposition 10); the latter has as basis the tensor product of the bases \((1, e_i)\) of the \(\bigwedge(M_i)\) (II, § 3, no. 7, Corollary 2 to Proposition 7); thus we obtain all the elements

\[
u _ {1} \otimes u _ {2} \otimes \dots \otimes u _ {m}
\]

where either \(u_{i}=1\) or \(u_{i}=e_{i}\) ; if J is the set of indices i such that \(u_{i}=e_{i}\) , \(u_{1}\otimes u_{2}\otimes\cdots\otimes u_{m}\) is identical with \(e_{J}\) , which completes the proof.

COROLLARY 1. Suppose that \(\mathbf{L} = [1, m]\) ; then the basis \((e_{\mathrm{J}})_{\mathrm{J} \in \mathfrak{P}(\mathrm{L})}\) of \(\bigwedge (\mathbf{M})\) has \(2^{m}\) elements. For \(p > m\) , \(\bigwedge^{p}(\mathbf{M}) = \{0\}\) ; \(\bigwedge^{m}(\mathbf{M})\) has a basis consisting of a single element \(e_{\mathrm{L}}\) ; for \(0 \leqslant p \leqslant m\) the number of elements in the basis \((e_{\mathrm{J}})\) of \(\bigwedge^{p}(\mathbf{M})\) consisting of the \(e_{\mathrm{J}}\) such that \(\operatorname{Card}(\mathbf{J}) = p\) is

\[
\binom {m} {p} = \frac {m !}{p ! (m - p) !}.
\]

This follows from Set Theory, III, § 3, no. 5, Proposition 12 and Set Theory, III, § 5, no. 8, Corollary 1 to Proposition 11.

We return to the case where the set L in Theorem 1 is arbitrary and give explicitly the multiplication table (§ 1, no. 7) of the basis \((e_{J})\) . Given two finite subsets \(J\) , \(K\) of the totally ordered set \(L\) , we write

\[
\left\{ \begin{array}{l l} \rho_ {J, K} = 0 & \text { if } J \cap K \neq \varnothing \\ \rho_ {J, K} = (- 1) ^ {\nu} & \text { if } J \cap K = \varnothing \end{array} \right.\tag{19}
\]

where v denotes the number of ordered pairs \((\lambda, \mu) \in J \times K\) such that \(\lambda > \mu\) . Then Corollary 1 to Proposition 4 of no. 2 immediately implies the relation

\[
e _ {J} \wedge e _ {K} = \rho_ {J, K} e _ {J \cup K}.\tag{20}
\]

Note the formula

\[
\rho_ {J, K} \rho_ {K, J} = (- 1) ^ {j k}\tag{21}
\]

when \(\mathbf{J} \cap \mathbf{K} = \varnothing, j = \operatorname{Card}(\mathbf{J}), k = \operatorname{Card}(\mathbf{K})\) (no. 3, Corollary 2 to Proposition 5.)

COROLLARY 2. If M is a projective A-module, \(\bigwedge (\mathbf{M})\) is a projective A-module.

The proof is the same as that of § 5, no. 5, Corollary to Theorem 1 replacing T by \(\Lambda\) .

COROLLARY 3. Let M be a projective A-module and \((x_{i})_{1\leqslant i\leqslant n}\) a finite sequence of elements of M. For there to exist on M an alternating \(n\) -linear form \(f\) such that

\[
f (x _ {1}, x _ {2}, \dots , x _ {n}) \neq 0,
\]

it is necessary and sufficient that \(x_{1} \wedge x_{2} \wedge \cdots \wedge x_{n} \neq 0\) .

We know already (with no hypothesis on M) that the condition is necessary (no. 4, Proposition 7). Suppose now that M is projective and that

\[
x _ {1} \wedge x _ {2} \wedge \dots \wedge x _ {n} \neq 0.
\]

Then \(\bigwedge^n (\mathbf{M})\) is projective (Corollary 2) and hence the canonical mapping

\[
\bigwedge^ {n} (\mathbf {M}) \rightarrow (\bigwedge^ {n} (\mathbf {M})) ^ {* *}
\]

is injective (II, §2, no. 7, Corollary 4 to Proposition 13); we conclude that there exists a linear form \(g: \bigwedge^{n}(\mathbf{M}) \to \mathbf{A}\) such that \(g(x_{1} \wedge x_{2} \wedge \dots \wedge x_{n}) \neq 0\) . If \(f\) is the alternating \(n\) -linear form corresponding to \(g\) (no. 4, Proposition 7), then \(f(x_{1}, \ldots, x_{n}) \neq 0\) .

